

\begin{register}{H}{GPIO\_OUT\_REG}{0x0004}\label{regdesc:GPIOOUTREG}
\regfield{GPIO\_OUT\_DATA\_ORIG}{32}{0}{{0x000000}}%
\reglabel{Reset}\regnewline%
\begin{regdesc}\begin{reglist}
\label{fielddesc:GPIOOUTDATAORIG}\item [GPIO\_OUT\_DATA\_ORIG] 配置简单 GPIO 输出模式下 GPIO0 \verb+~+ GPIO27 的输出值。 \\
0：低电平\\
1：高电平\\
bit0 \verb+~+ bit27 分别对应 GPIO0 \verb+~+ GPIO27 的输出值。bit28 \verb+~+ bit31 无效。\\
(R/W/SC/WTC)
\end{reglist}\end{regdesc}
\end{register}


\begin{register}{H}{GPIO\_OUT\_W1TS\_REG}{0x0008}\label{regdesc:GPIOOUTW1TSREG}
\regfield{GPIO\_OUT\_W1TS}{32}{0}{{0x000000}}%
\reglabel{Reset}\regnewline%
\begin{regdesc}\begin{reglist}
\label{fielddesc:GPIOOUTW1TS}\item [GPIO\_OUT\_W1TS] 配置是否置位 GPIO0 \verb+~+ GPIO27 输出寄存器 \hyperref[regdesc:GPIOOUTREG]{GPIO\_OUT\_REG}。\\
0：不置位\\
1：\hyperref[regdesc:GPIOOUTREG]{GPIO\_OUT\_REG} 中相应位也将置 1\\
bit0 \verb+~+ bit27 分别对应 GPIO0 \verb+~+ GPIO27。bit28 \verb+~+ bit31 无效。
推荐使用此寄存器来置位 \hyperref[regdesc:GPIOOUTREG]{GPIO\_OUT\_REG}。\\
(WT)
\end{reglist}\end{regdesc}
\end{register}


\begin{register}{H}{GPIO\_OUT\_W1TC\_REG}{0x000C}\label{regdesc:GPIOOUTW1TCREG}
\regfield{GPIO\_OUT\_W1TC}{32}{0}{{0x000000}}%
\reglabel{Reset}\regnewline%
\begin{regdesc}\begin{reglist}
\label{fielddesc:GPIOOUTW1TC}\item [GPIO\_OUT\_W1TC] 配置是否清除 GPIO0 \verb+~+ GPIO27 输出寄存器 \hyperref[regdesc:GPIOOUTREG]{GPIO\_OUT\_REG}。\\
0：不清除\\
1：\hyperref[regdesc:GPIOOUTREG]{GPIO\_OUT\_REG} 中相应位将被清除\\
bit0 \verb+~+ bit27 分别对应 GPIO0 \verb+~+ GPIO27。bit28 \verb+~+ bit31 无效。
推荐使用该寄存器来清除 \hyperref[regdesc:GPIOOUTREG]{GPIO\_OUT\_REG}。\\
(WT)
\end{reglist}\end{regdesc}
\end{register}


\begin{register}{H}{GPIO\_ENABLE\_REG}{0x0020}\label{regdesc:GPIOENABLEREG}
\regfield{GPIO\_ENABLE\_DATA}{32}{0}{{0x000000}}%
\reglabel{Reset}\regnewline%
\begin{regdesc}\begin{reglist}
\label{fielddesc:GPIOENABLEDATA}\item [GPIO\_ENABLE\_DATA] 配置是否使能 GPIO0 \verb+~+ GPIO27 输出。\\
0：不使能\\
1：使能\\
bit0 \verb+~+ bit27 分别对应 GPIO0 \verb+~+ GPIO27。bit28 \verb+~+ bit31 无效。\\
(R/W/WTC)
\end{reglist}\end{regdesc}
\end{register}


\begin{register}{H}{GPIO\_ENABLE\_W1TS\_REG}{0x0024}\label{regdesc:GPIOENABLEW1TSREG}
\regfield{GPIO\_ENABLE\_W1TS}{32}{0}{{0x000000}}%
\reglabel{Reset}\regnewline%
\begin{regdesc}\begin{reglist}
\label{fielddesc:GPIOENABLEW1TS}\item [GPIO\_ENABLE\_W1TS] 配置是否使能 GPIO0 \verb+~+ GPIO27 的输出使能寄存器 \hyperref[regdesc:GPIOENABLEREG]{GPIO\_ENABLE\_REG}。\\
0：不使能\\
1：\hyperref[regdesc:GPIOENABLEREG]{GPIO\_ENABLE\_REG} 中相应位也将置 1\\
bit0 \verb+~+ bit27 分别对应 GPIO0 \verb+~+ GPIO27。bit28 \verb+~+ bit31 无效。
推荐使用该寄存器来置位 \hyperref[regdesc:GPIOENABLEREG]{GPIO\_ENABLE\_REG}。\\
(WT)
\end{reglist}\end{regdesc}
\end{register}


\begin{register}{H}{GPIO\_ENABLE\_W1TC\_REG}{0x0028}\label{regdesc:GPIOENABLEW1TCREG}
\regfield{GPIO\_ENABLE\_W1TC}{32}{0}{{0x000000}}%
\reglabel{Reset}\regnewline%
\begin{regdesc}\begin{reglist}
\label{fielddesc:GPIOENABLEW1TC}\item [GPIO\_ENABLE\_W1TC] 配置是否清除 GPIO0 \verb+~+ GPIO27 输出使能寄存器 \hyperref[regdesc:GPIOENABLEREG]{GPIO\_ENABLE\_REG}。 \\
0：不清除\\
1：\hyperref[regdesc:GPIOENABLEREG]{GPIO\_ENABLE\_REG} 中相应位将被清除\\
bit0 \verb+~+ bit27 分别对应 GPIO0 \verb+~+ 27。bit28 \verb+~+ bit31 无效。
推荐使用此寄存器清除 \hyperref[regdesc:GPIOENABLEREG]{GPIO\_ENABLE\_REG}。\\
(WT)
\end{reglist}\end{regdesc}
\end{register}

\begin{register}{H}{GPIO\_STRAP\_REG}{0x0038}\label{regdesc:GPIOSTRAPREG}
\regfield{(reserved)}{16}{16}{0000000000000000}%
\regfield{GPIO\_STRAPPING}{16}{0}{{0x00}}%
\reglabel{Reset}\regnewline%
\begin{regdesc}\begin{reglist}
\label{fielddesc:GPIOSTRAPPING}\item [GPIO\_STRAPPING] 表示 GPIO Strapping 管脚的值。
\begin{itemize}
    \item bit0：表示 GPIO2 管脚的值（仅在使用 SPI Download Boot 模式时，GPIO2 才用作 Strapping 管脚）
    \item bit1：表示 GPIO3 管脚的值（仅在使用 SPI Download Boot 模式时，GPIO3 才用作 Strapping 管脚）
    \item bit2：表示 GPIO8 管脚的值（该值同时会受到 \hyperref[fielddesc:EFUSEDISFORCEDOWNLOAD]{EFUSE\_DIS\_FORCE\_DOWNLOAD} 和 LP\_AON\_FORCE\_DOWNLOAD\_BOOT 的影响）
    \item bit3：表示 GPIO9 管脚的值（该值同时会受到 \hyperref[fielddesc:EFUSEDISFORCEDOWNLOAD]{EFUSE\_DIS\_FORCE\_DOWNLOAD} 和 LP\_AON\_FORCE\_DOWNLOAD\_BOOT 的影响）
    \item bit4：表示 GPIO25 管脚的值
    \item bit5 \verb+~+ bit15：无效
\end{itemize}
关于 GPIO Strapping 管脚详细信息，请参考 \ref{mod:bootctrl} \textit{\nameref{mod:bootctrl}} 章节。\\
(RO)
\end{reglist}\end{regdesc}
\end{register}


\begin{register}{H}{GPIO\_IN\_REG}{0x003C}\label{regdesc:GPIOINREG}
\regfield{GPIO\_IN\_DATA\_NEXT}{32}{0}{{0x000000}}%
\reglabel{Reset}\regnewline%
\begin{regdesc}\begin{reglist}
\label{fielddesc:GPIOINDATANEXT}\item [GPIO\_IN\_DATA\_NEXT] 表示 GPIO0 \verb+~+ GPIO27 的输入值。每一位均代表一个管脚的输入值：\\
0：低电平 \\
1：高电平 \\
bit0 \verb+~+ bit27 分别对应 GPIO0 \verb+~+ GPIO27，bit28 \verb+~+ bit31 无效。\\
(RO)
\end{reglist}\end{regdesc}
\end{register}

\begin{register}{H}{GPIO\_STATUS\_REG}{0x0044}\label{regdesc:GPIOSTATUSREG}
\regfield{GPIO\_STATUS\_INTERRUPT}{32}{0}{{0x000000}}%
\reglabel{Reset}\regnewline%
\begin{regdesc}\begin{reglist}
\label{fielddesc:GPIOSTATUSINTERRUPT}\item [GPIO\_STATUS\_INTERRUPT] GPIO0 \verb+~+ GPIO27 的中断状态寄存器，并且软件可配置该寄存器的值。
\begin{itemize}
    \item bit0 \verb+~+ bit27 分别对应 GPIO0 \verb+~+ GPIO27。bit28 \verb+~+ bit31 无效。
    \item 每一位可查询对应 GPIO 的中断状态：
    \begin{itemize}
        \item 0：表示对应 GPIO 未产生 \hyperref[fielddesc:GPIOPINNINTTYPE]{GPIO\_PIN\regindex{n}\_INT\_TYPE} 所选择的中断，或软件配置该位寄存器的值为 0
        \item 1：表示对应 GPIO 产生了 \hyperref[fielddesc:GPIOPINNINTTYPE]{GPIO\_PIN\regindex{n}\_INT\_TYPE} 所选择的中断，或软件配置该位寄存器的值为 1
    \end{itemize}
\end{itemize}
(R/W/WTC)
\end{reglist}\end{regdesc}
\end{register}


\begin{register}{H}{GPIO\_STATUS\_W1TS\_REG}{0x0048}\label{regdesc:GPIOSTATUSW1TSREG}
\regfield{GPIO\_STATUS\_W1TS}{32}{0}{{0x000000}}%
\reglabel{Reset}\regnewline%
\begin{regdesc}\begin{reglist}
\label{fielddesc:GPIOSTATUSW1TS}\item [GPIO\_STATUS\_W1TS] 配置是否置位 GPIO0 \verb+~+ GPIO27 的中断状态寄存器 \hyperref[fielddesc:GPIOSTATUSINTERRUPT]{GPIO\_STATUS\_INTERRUPT}。
\begin{itemize}
    \item bit0 \verb+~+ bit27 分别对应 GPIO0 \verb+~+ GPIO27。bit28 \verb+~+ bit31 无效。
    \item 每一位置 1，则 \hyperref[fielddesc:GPIOSTATUSINTERRUPT]{GPIO\_STATUS\_INTERRUPT} 中相应位也将置 1。
    \item 推荐使用此寄存器置位 \hyperref[fielddesc:GPIOSTATUSINTERRUPT]{GPIO\_STATUS\_INTERRUPT}。
\end{itemize}
(WT)
\end{reglist}\end{regdesc}
\end{register}


\begin{register}{H}{GPIO\_STATUS\_W1TC\_REG}{0x004C}\label{regdesc:GPIOSTATUSW1TCREG}
\regfield{GPIO\_STATUS\_W1TC}{32}{0}{{0x000000}}%
\reglabel{Reset}\regnewline%
\begin{regdesc}\begin{reglist}
\label{fielddesc:GPIOSTATUSW1TC}\item [GPIO\_STATUS\_W1TC] 配置是否清除 GPIO0 \verb+~+ GPIO27 的中断状态寄存器 \hyperref[fielddesc:GPIOSTATUSINTERRUPT]{GPIO\_STATUS\_INTERRUPT}。
\begin{itemize}
    \item bit0 \verb+~+ bit27 分别对应 GPIO0 \verb+~+ GPIO27。bit28 \verb+~+ bit31 无效。
    \item 每一位置 1，则 \hyperref[fielddesc:GPIOSTATUSINTERRUPT]{GPIO\_STATUS\_INTERRUPT} 中相应位将被清除。
    \item 推荐使用此寄存器来清除 \hyperref[fielddesc:GPIOSTATUSINTERRUPT]{GPIO\_STATUS\_INTERRUPT}。
\end{itemize}(WT)
\end{reglist}\end{regdesc}
\end{register}

\begin{register}{H}{GPIO\_PCPU\_INT\_REG}{0x005C}\label{regdesc:GPIOPCPUINTREG}
\regfield{GPIO\_PROCPU\_INT}{32}{0}{{0x000000}}%
\reglabel{Reset}\regnewline%
\begin{regdesc}\begin{reglist}
\label{fielddesc:GPIOPROCPUINT}\item [GPIO\_PROCPU\_INT] 表示 GPIO0 \verb+~+ GPIO27 CPU 中断状态。
\begin{itemize}
    \item bit0 \verb+~+ bit27 分别对应 GPIO0 \verb+~+ GPIO27。bit28 \verb+~+ bit31 无效。
    \item 如果 \hyperref[regdesc:GPIOPINNREG]{GPIO\_PIN\regindex{n}\_REG} 中 bit13 有效，即使能 CPU 中断，则此寄存器所示的中断状态应与 \hyperref[regdesc:GPIOSTATUSREG]{GPIO\_STATUS\_REG} 中相应位的中断状态一致，否则为 0。
    \begin{itemize}
        \item 0：表示未使能 CPU 中断，或对应 GPIO 未产生 \hyperref[fielddesc:GPIOPINNINTTYPE]{GPIO\_PIN\regindex{n}\_INT\_TYPE} 所选择的中断
        \item 1：表示使能 CPU 中断后，对应 GPIO 产生了 \hyperref[fielddesc:GPIOPINNINTTYPE]{GPIO\_PIN\regindex{n}\_INT\_TYPE} 所选择的中断
    \end{itemize}
\end{itemize}
(RO)
\end{reglist}\end{regdesc}
\end{register}





\begin{register}{H}{GPIO\_PIN\regindex{n}\_REG (\regindex{n}: 0-27)}{0x0074+4*\regindex{n}}\label{regdesc:GPIOPINNREG}
\regfield{(reserved)}{14}{18}{00000000000000}%
\regfield{GPIO\_PIN\regindex{n}\_INT\_ENA}{5}{13}{{0x0}}%
\regfield{(reserved)}{2}{11}{{0x0}}%
\regfield{GPIO\_PIN\regindex{n}\_WAKEUP\_ENABLE}{1}{10}{{0}}%
\regfield{GPIO\_PIN\regindex{n}\_INT\_TYPE}{3}{7}{{0x0}}%
\regfield{(reserved)}{2}{5}{00}%
\regfield{GPIO\_PIN\regindex{n}\_SYNC1\_BYPASS}{2}{3}{{0x0}}%
\regfield{GPIO\_PIN\regindex{n}\_PAD\_DRIVER}{1}{2}{{0}}%
\regfield{GPIO\_PIN\regindex{n}\_SYNC2\_BYPASS}{2}{0}{{0x0}}%
\reglabel{Reset}\regnewline%
\begin{regdesc}\begin{reglist}
\label{fielddesc:GPIOPINNSYNC2BYPASS}\item [GPIO\_PIN\regindex{n}\_SYNC2\_BYPASS] 配置是否使能 GPIO 输入数据在第二级同步中 IO MUX 运行时钟上升沿同步或下降沿同步。\\
0：不进行同步\\
1：在下降沿进行同步\\
2：在上升沿进行同步\\
3：在上升沿进行同步\\
(R/W)
\label{fielddesc:GPIOPINNPADDRIVER}\item [GPIO\_PIN\regindex{n}\_PAD\_DRIVER] 配置选择管脚的输出驱动模式。\\
0：正常输出\\
1：开漏输出 \\(R/W)
\label{fielddesc:GPIOPINNSYNC1BYPASS}\item [GPIO\_PIN\regindex{n}\_SYNC1\_BYPASS] 配置是否使能 GPIO 输入数据在第一级同步中 IO MUX 运行时钟上升沿同步或下降沿同步。\\
0：不进行同步\\
1：在下降沿进行同步\\
2：在上升沿进行同步\\
3：在上升沿进行同步\\
(R/W)
\label{fielddesc:GPIOPINNINTTYPE}\item [GPIO\_PIN\regindex{n}\_INT\_TYPE] 配置 GPIO 中断类型。\\
0：关闭 GPIO 中断\\
1：上升沿触发\\
2：下降沿触发\\
3：任一沿触发\\
4：低电平触发\\
5：高电平触发\\ (R/W)
\label{fielddesc:GPIOPINNWAKEUPENABLE}\item [GPIO\_PIN\regindex{n}\_WAKEUP\_ENABLE] 配置是否使能 GPIO 唤醒功能。\\
0：不使能\\
1：使能\\
该功能仅能将 CPU 从 Light-sleep 模式唤醒。 \\(R/W)
\item [见下页]
\end{reglist}\end{regdesc}
\end{register}

\addtocounter{Regfloat}{-1}
\begin{register}{H}{GPIO\_PIN\regindex{n}\_REG (\regindex{n}: 0-27)}{0x0074+4*\regindex{n}}
\begin{regdesc}\begin{reglist}
\item [接上页]
\label{fielddesc:GPIOPINNINTENA}\item [GPIO\_PIN\regindex{n}\_INT\_ENA] 配置是否使能 CPU 中断或 CPU 非屏蔽中断。
\begin{itemize}
    \item bit13：配置是否使能 CPU 中断：\\
    0：不使能\\
    1：使能\\
    \item bit14：配置是否使能 CPU 非屏蔽中断： \\
    0：不使能\\
    1：使能\\
    \item bit15 \verb+~+ bit17：无效
    \end{itemize} (R/W)
\end{reglist}\end{regdesc}
\end{register}


\begin{register}{H}{GPIO\_STATUS\_NEXT\_REG}{0x014C}\label{regdesc:GPIOSTATUSNEXTREG}
\regfield{GPIO\_STATUS\_INTERRUPT\_NEXT}{32}{0}{{0x000000}}%
\reglabel{Reset}\regnewline%
\begin{regdesc}\begin{reglist}
\label{fielddesc:GPIOSTATUSINTERRUPTNEXT}\item [GPIO\_STATUS\_INTERRUPT\_NEXT] 表示 GPIO0 \verb+~+ GPIO27 的中断源信号状态。\\
bit0 \verb+~+ bit27 分别对应 GPIO0 \verb+~+ 27。bit28 \verb+~+ bit31 无效，每个 bit 表示：
\begin{itemize}
        \item 0：表示对应 GPIO 未产生 \hyperref[fielddesc:GPIOPINNINTTYPE]{GPIO\_PIN\regindex{n}\_INT\_TYPE} 所选择的中断
        \item 1：表示对应 GPIO 产生了 \hyperref[fielddesc:GPIOPINNINTTYPE]{GPIO\_PIN\regindex{n}\_INT\_TYPE} 所选择的中断
\end{itemize}
上述中断可以为上升沿中断、下降沿中断、电平敏感中断或任一沿中断。\\  (RO)

\end{reglist}\end{regdesc}
\end{register}


\begin{register}{H}{GPIO\_FUNC\regindex{n}\_IN\_SEL\_CFG\_REG (\regindex{n}: 0-127)}{0x0154+4*\regindex{n}}\label{regdesc:GPIOFUNCNINSELCFGREG}
\regfield{(reserved)}{24}{8}{000000000000000000000000}%
\regfield{GPIO\_SIG\regindex{n}\_IN\_SEL}{1}{7}{{0}}%
\regfield{GPIO\_FUNC\regindex{n}\_IN\_INV\_SEL}{1}{6}{{0}}%
\regfield{GPIO\_FUNC\regindex{n}\_IN\_SEL}{6}{0}{{0x0}}%
\reglabel{Reset}\regnewline%
\begin{regdesc}\begin{reglist}
\label{fielddesc:GPIOFUNCNINSEL}\item [GPIO\_FUNC\regindex{n}\_IN\_SEL] 配置从 28 个 GPIO 中选择一个管脚连接输入信号 \regindex{n}。\\
0：选择 GPIO0\\
1：选择 GPIO1\\
......\\
26：选择 GPIO26\\
27：选择 GPIO27\\
或者\\
0x38：选择恒高电平输入\\
0x3C：选择恒低电平输入\\
(R/W)
\label{fielddesc:GPIOFUNCNININVSEL}\item [GPIO\_FUNC\regindex{n}\_IN\_INV\_SEL] 配置是否反转输入值。\\
0：不反转\\
1：反转\\ (R/W)
\label{fielddesc:GPIOSIGNINSEL}\item [GPIO\_SIG\regindex{n}\_IN\_SEL] 配置是否通过 GPIO 交换矩阵连接信号。\\
0：旁路 GPIO 交换矩阵，即直接通过 IO MUX 连接信号与外设\\
1：通过 GPIO 交换矩阵连接信号与外设\\  (R/W)
\end{reglist}\end{regdesc}
\end{register}


\begin{register}{H}{GPIO\_FUNC\regindex{n}\_OUT\_SEL\_CFG\_REG (\regindex{n}: 0-27)}{0x0554+4*\regindex{n}}\label{regdesc:GPIOFUNCNOUTSELCFGREG}
\regfield{(reserved)}{21}{11}{000000000000000000000}%
\regfield{GPIO\_FUNC0\_OEN\_INV\_SEL}{1}{10}{{0}}%
\regfield{GPIO\_FUNC0\_OEN\_SEL}{1}{9}{{0}}%
\regfield{GPIO\_FUNC0\_OUT\_INV\_SEL}{1}{8}{{0}}%
\regfield{GPIO\_FUNC0\_OUT\_SEL}{8}{0}{{0x80}}%
\reglabel{Reset}\regnewline%
\begin{regdesc}\begin{reglist}
\label{fielddesc:GPIOFUNCNOUTSEL}\item [GPIO\_FUNC\regindex{n}\_OUT\_SEL] 配置从 128 个外设信号中选择一个信号 \regindex{Y} (0 <= \regindex{Y} < 128) 输出至 GPIO\regindex{n}。\\
0：选择信号 0\\
1：选择信号 1\\
......\\
126：选择信号 126\\
127：选择信号 127\\
或者\\
128：\hyperref[regdesc:GPIOOUTREG]{GPIO\_OUT\_REG} 和 \hyperref[regdesc:GPIOENABLEREG]{GPIO\_ENABLE\_REG} 中的 bit\regindex{n} 将用作输出值和输出使能信号。

信号列表见表 \ref{tab:iomuxgpio-periph-signals-via-gpio-matrix}。

(R/W/SC)
\label{fielddesc:GPIOFUNCNOUTINVSEL}\item [GPIO\_FUNC\regindex{n}\_OUT\_INV\_SEL] 配置是否反转输出值。\\
0：不反转\\
1：反转\\ (R/W/SC)
\label{fielddesc:GPIOFUNCNOENSEL}\item [GPIO\_FUNC\regindex{n}\_OEN\_SEL] 配置选择输出使能信号。\\
0：使用外设的输出使能信号\\
1：强制使用 \hyperref[regdesc:GPIOENABLEREG]{GPIO\_ENABLE\_REG} 的 bit\regindex{n} 用作输出使能信号 \\(R/W)
\label{fielddesc:GPIOFUNCNOENINVSEL}\item [GPIO\_FUNC\regindex{n}\_OEN\_INV\_SEL] 配置是否反转输出使能信号。\\
0：不反转\\
1：反转\\ (R/W)
\end{reglist}\end{regdesc}
\end{register}


\begin{register}{H}{GPIO\_CLOCK\_GATE\_REG}{0x062C}\label{regdesc:GPIOCLOCKGATEREG}
\regfield{(reserved)}{31}{1}{0000000000000000000000000000000}%
\regfield{GPIO\_CLK\_EN}{1}{0}{{1}}%
\reglabel{Reset}\regnewline%
\begin{regdesc}\begin{reglist}
\label{fielddesc:GPIOCLKEN}\item [GPIO\_CLK\_EN] 配置是否使能时钟门控。\\
0：不使能\\
1：使能，则时钟自由运转。\\(R/W)
\end{reglist}\end{regdesc}
\end{register}


\begin{register}{H}{GPIO\_DATE\_REG}{0x06FC}\label{regdesc:GPIOREGDATEREG}
\regfield{(reserved)}{4}{28}{0000}%
\regfield{GPIO\_DATE}{28}{0}{{0x2201120}}%
\reglabel{Reset}\regnewline%
\begin{regdesc}\begin{reglist}
\label{fielddesc:GPIOREGDATE}\item [GPIO\_DATE] 版本控制寄存器。 (R/W)
\end{reglist}\end{regdesc}
\end{register}
